% CONCLUSION. Numbers ONLY via \R... macros.
We asked whether measuring how much each type of context contributes to task success can be reused for budgeted
context assembly. In these experiments, measured context value was more useful for diagnosing which kinds of context
matter than for ranking individual blocks. Embedding relevance remained difficult to beat at the instance level: neither
the MCV compiler nor a hybrid built on the profile improved on it, and supervision from single-block deletions was too
sparse to learn block value. Type-level ablations nevertheless exposed concentrated, redundant and low-value context,
and the largest token savings came from deciding when to stop sending context. These results motivate context-assembly
designs that keep type-level valuation and block-level relevance separate, and we release every log so that the
negative results can be checked and extended.

\textbf{Use of AI assistance.} An AI coding assistant was used to write and test code, run the experiments, and draft
and edit text. All numbers in the paper are generated from the released logs by scripts, and every reference was
verified against its DOI record. The author reviewed the content and is responsible for it.

\textbf{Artifact availability.} The \texttt{mcv} package, data adapters (with the licenses of HotpotQA and BFCL),
pre-registration, raw logs of every model call, and the scripts that regenerate all figures, tables and numbers with one
command will be released under an open license with the preprint.
